% ============================================================
\chapter{Transactions and Concurrency Control}
\label{ch:transactions}
% ============================================================

Imagine a bank transfer: one account is debited and another is credited. What happens if the system crashes between the two operations? The money has vanished. It was to prevent such catastrophes that Jim Gray, in the 1970s at IBM, formalized the notion of a \emph{transaction} and the ACID properties (Atomicity, Consistency, Isolation, Durability). A transaction is a logical ``all or nothing'': either all its operations succeed, or none takes effect. When multiple transactions execute concurrently, concurrency control (locking, timestamping, MVCC) guarantees that the result is equivalent to a serial execution. Gray received the Turing Award in 1998 for this foundational work.

\begin{intuition}[title=\textbf{Central idea}]
A transaction is a logical unit of work that must be executed atomically,
consistently, in isolation, and durably (ACID). Concurrency control ensures
that the simultaneous execution of multiple transactions produces a correct
result, as if they had been executed sequentially.
\end{intuition}

% ------------------------------------------------------------
\section{The concept of a transaction}
% ------------------------------------------------------------

\begin{definition}[Transaction]
\index{transaction}
A \textbf{transaction} is a sequence of read and write operations on the
database, delimited by a \texttt{BEGIN} and terminated by a \texttt{COMMIT}
(validation) or a \texttt{ROLLBACK} (abort).
\end{definition}

\begin{minted}{python}
# Example in Python with psycopg2
import psycopg2

conn = psycopg2.connect("dbname=bank")
cur = conn.cursor()
try:
    cur.execute("UPDATE accounts SET balance = balance - 100 WHERE id = 1")
    cur.execute("UPDATE accounts SET balance = balance + 100 WHERE id = 2")
    conn.commit()   # Both operations are committed together
except Exception:
    conn.rollback()  # No modification is applied
finally:
    cur.close()
    conn.close()
\end{minted}

% ------------------------------------------------------------
\section{ACID properties}
% ------------------------------------------------------------

\begin{definition}[ACID properties]
\index{ACID}
The four fundamental properties of a transaction are:
\begin{enumerate}
  \item \textbf{Atomicity} --- All operations are executed or none.
  \item \textbf{Consistency} --- The transaction brings the database from
        one consistent state to another.
  \item \textbf{Isolation} --- Concurrent transactions do not see each
        other's intermediate modifications.
  \item \textbf{Durability} --- Once committed, modifications survive any
        crash.
\end{enumerate}
\end{definition}

\begin{remark}
Atomicity and durability are ensured by the \emph{Write-Ahead Log} (WAL).
Consistency is ensured by integrity constraints. Isolation is ensured by
concurrency control mechanisms.
\end{remark}

% ------------------------------------------------------------
\section{Concurrency anomalies}
% ------------------------------------------------------------

Without control, concurrent execution can produce anomalies:

\begin{definition}[Concurrency anomalies]
\index{concurrency anomalies}
\begin{itemize}
  \item \textbf{Dirty read} --- $T_2$ reads data modified by $T_1$ that has
        not yet committed.
  \item \textbf{Non-repeatable read} --- $T_1$ reads the same data twice and
        gets different values because $T_2$ modified it in between.
  \item \textbf{Phantom read} --- $T_1$ executes a query twice and gets
        different rows because $T_2$ inserted or deleted rows.
  \item \textbf{Lost update} --- Two transactions read the same value,
        modify it, and write it back: one modification is lost.
\end{itemize}
\end{definition}

% ------------------------------------------------------------
\section{Isolation levels}
% ------------------------------------------------------------

\begin{definition}[SQL isolation levels]
\index{isolation levels}
The SQL standard defines four isolation levels:
\begin{center}
\begin{tabular}{lccc}
\hline
\textbf{Level} & \textbf{Dirty read} & \textbf{Non-repeatable} & \textbf{Phantom} \\
\hline
\texttt{READ UNCOMMITTED} & Possible & Possible & Possible \\
\texttt{READ COMMITTED}   & No       & Possible & Possible \\
\texttt{REPEATABLE READ}  & No       & No       & Possible \\
\texttt{SERIALIZABLE}     & No       & No       & No \\
\hline
\end{tabular}
\end{center}
\end{definition}

\begin{attention}
The default level varies by DBMS: \texttt{READ COMMITTED} for PostgreSQL
and Oracle, \texttt{REPEATABLE READ} for MySQL/InnoDB. Choosing too weak a
level exposes you to anomalies; too strong a level reduces performance.
\end{attention}

\begin{minted}{python}
-- Setting the isolation level in SQL
SET TRANSACTION ISOLATION LEVEL SERIALIZABLE;
BEGIN;
    SELECT balance FROM accounts WHERE id = 1;
    UPDATE accounts SET balance = balance - 100 WHERE id = 1;
COMMIT;
\end{minted}

% ------------------------------------------------------------
\section{Serializability theory}
% ------------------------------------------------------------

\begin{definition}[Schedule]
\index{schedule}
A \textbf{schedule} $S$ is a sequence of operations (reads, writes,
commits, rollbacks) from multiple transactions that preserves the internal
order of each transaction.
\end{definition}

\begin{definition}[Conflict serializability]
\index{serializability}
Two operations \textbf{conflict} if they operate on the same data item and
at least one is a write. A schedule $S$ is \textbf{conflict-serializable}
if it is conflict-equivalent to some serial schedule.
\end{definition}

\begin{theorem}[Precedence graph test]
\index{precedence graph}
Construct the \textbf{precedence graph} $G = (V, E)$ where $V$ is the set
of transactions and $(T_i, T_j) \in E$ if an operation of $T_i$ conflicts
with an operation of $T_j$ and precedes it. The schedule is conflict-serializable
if and only if $G$ is acyclic.
\end{theorem}

\begin{example}
Let $S : r_1(A),\, w_1(A),\, r_2(A),\, w_2(A),\, r_1(B),\, w_1(B),\, r_2(B),\, w_2(B)$.
Conflicts: $w_1(A) \to r_2(A)$, $w_1(A) \to w_2(A)$,
$w_1(B) \to r_2(B)$, $w_1(B) \to w_2(B)$. The graph has a single edge
$T_1 \to T_2$, so $S$ is serializable (equivalent to $T_1, T_2$).
\end{example}

% ------------------------------------------------------------
\section{Two-phase locking (2PL)}
% ------------------------------------------------------------

\begin{definition}[2PL protocol]
\index{two-phase locking}
The \textbf{two-phase locking} (2PL) protocol requires each transaction
to go through two phases:
\begin{enumerate}
  \item \textbf{Growing phase}: the transaction acquires locks (shared for
        reads, exclusive for writes) without releasing any.
  \item \textbf{Shrinking phase}: the transaction releases locks without
        acquiring new ones.
\end{enumerate}
\end{definition}

\begin{theorem}[Correctness of 2PL]
Every schedule produced by the 2PL protocol is conflict-serializable.
\end{theorem}

\begin{definition}[Strict and rigorous 2PL]
\begin{itemize}
  \item \textbf{Strict 2PL}: exclusive locks are released only at
        \texttt{COMMIT}/\texttt{ROLLBACK}. Avoids cascading aborts.
  \item \textbf{Rigorous 2PL}: \emph{all} locks are released at
        \texttt{COMMIT}/\texttt{ROLLBACK}. Most commonly used mode.
\end{itemize}
\end{definition}

% ------------------------------------------------------------
\section{Deadlock detection and prevention}
% ------------------------------------------------------------

\begin{definition}[Deadlock]
\index{deadlock}
A \textbf{deadlock} occurs when two or more transactions wait for each
other, each holding a lock requested by the other.
\end{definition}

\begin{proposition}[Wait-for graph detection]
Construct the \textbf{wait-for graph}: an edge $(T_i, T_j)$ means $T_i$
is waiting for a lock held by $T_j$. A deadlock exists if and only if the
graph contains a cycle.
\end{proposition}

Resolution strategies:
\begin{itemize}
  \item \textbf{Periodic detection}: check for cycles and abort the least
        costly transaction (victim selection).
  \item \textbf{Timestamp-based prevention}: Wait-Die or Wound-Wait schemes.
  \item \textbf{Timeout}: abort any transaction that waits too long.
\end{itemize}

% ------------------------------------------------------------
\section{Multi-Version Concurrency Control (MVCC)}
% ------------------------------------------------------------

\begin{definition}[MVCC]
\index{MVCC}
\textbf{MVCC} (Multi-Version Concurrency Control) maintains multiple
versions of each row. Reads never block writes and vice versa:
\begin{itemize}
  \item Each transaction sees a consistent snapshot of the database at its
        start time.
  \item Writes create new versions without overwriting old ones.
  \item A garbage collector (vacuum) removes obsolete versions.
\end{itemize}
\end{definition}

\begin{remark}
PostgreSQL uses MVCC for all isolation levels. MySQL/InnoDB uses MVCC for
\texttt{READ COMMITTED} and \texttt{REPEATABLE READ}, adding additional
locks for \texttt{SERIALIZABLE}.
\end{remark}

\begin{example}
Under MVCC, transaction $T_1$ can read version $v_1$ of a row while $T_2$
writes version $v_2$: no locking is needed for the read.
\begin{minted}{python}
-- PostgreSQL: observing MVCC
BEGIN;
SELECT xmin, xmax, * FROM accounts WHERE id = 1;
-- xmin = ID of the transaction that created this version
-- xmax = ID of the transaction that deleted it (0 if active)
COMMIT;
\end{minted}
\end{example}

% ------------------------------------------------------------
\section{Write-Ahead Logging (WAL)}
% ------------------------------------------------------------

\begin{definition}[Write-Ahead Logging]
\index{WAL}
The \textbf{WAL} (Write-Ahead Log) is a sequential log where every
modification is recorded \emph{before} being applied to the data pages.
In case of a crash, the log enables:
\begin{itemize}
  \item \textbf{REDO}: replay modifications of committed transactions.
  \item \textbf{UNDO}: reverse modifications of uncommitted transactions.
\end{itemize}
\end{definition}

% ------------------------------------------------------------
\section{Key formulas}
% ------------------------------------------------------------

\begin{keyformulas}
\begin{itemize}
  \item \textbf{ACID}: Atomicity, Consistency, Isolation, Durability
  \item \textbf{Serializability}: acyclic precedence graph $\Leftrightarrow$ serializable
  \item \textbf{2PL}: growing phase $\to$ lock point $\to$ shrinking phase
  \item \textbf{Deadlock}: cycle in the wait-for graph
  \item \textbf{MVCC}: non-blocking reads via multi-version snapshots
\end{itemize}
\end{keyformulas}

% ------------------------------------------------------------
\section{Exercises}
% ------------------------------------------------------------

\begin{exercise}
Given the schedule $S : r_1(A),\, r_2(B),\, w_1(B),\, w_2(A)$, construct
the precedence graph. Is $S$ serializable?
\end{exercise}

\begin{exercise}
Give an example of a deadlock between two transactions operating on two
bank accounts. Propose a solution using a fixed lock acquisition order.
\end{exercise}

\begin{exercise}
Write a Python script demonstrating a phantom read under \texttt{READ
COMMITTED} in PostgreSQL, then show that it disappears under
\texttt{SERIALIZABLE}.
\end{exercise}

\begin{exercise}
Explain why strict 2PL prevents cascading aborts. Give a counterexample
with basic 2PL.
\end{exercise}

\begin{exercise}
In an MVCC system, transaction $T_1$ reads row $r$ (version $v_1$), then
$T_2$ modifies $r$ and commits (version $v_2$), then $T_1$ re-reads $r$.
Which version does $T_1$ see under \texttt{READ COMMITTED}? Under
\texttt{REPEATABLE READ}? Justify.
\end{exercise}
